\section{Introduction}
\label{sec:introduction}

Modern cryptographic systems are increasingly used to protect privacy, verify computation, and support trustless applications \cite{parno2016pinocchio, rauzy2014countermeasures}. As these systems grow more complex, their security depends not only on the underlying mathematical protocol, but also on the correctness of the software that implements it. A small implementation mistake may cause a program to accept an invalid proof, reject a valid proof, mishandle an input value, or enforce a different statement from the one intended by the developer \cite{wen2024practical, tang2024zero}. These errors are especially dangerous because the program may still execute normally and produce no crash or memory-safety failure.

Zero-knowledge proof systems are a representative example of this challenge. A zero-knowledge proof program usually contains multiple components, including circuit construction, proof generation, and proof verification \cite{groth2016size}. Each component must interpret the public inputs, private witness values, cryptographic keys, and proof objects consistently. If any component applies an incorrect encoding rule, omits a necessary constraint, or mishandles a boundary condition, the final verification result may no longer reflect the intended statement. Such errors are cryptographic logic errors rather than ordinary software bugs.

There are several protocol families in Zero-knowledge
proof systems, such as zkSNARKs, zkSTARKs, and Bulletproofs. Despite their strong theoretical security guarantees, the complexity of these systems makes implementation and deployment errors a serious practical concern. A representative example is Aztec 2.0, which used a PLONK-family zkSNARK for private asset transfers on Ethereum \cite{aztec_vuln_2021}. A previous vulnerability in its recursive proof verification logic allowed fake join-split proofs, including double-spend transactions, to pass without being detected by either the rollup circuit logic or the verifier smart contract logic, thereby allowing the same asset to be spent more than once. Additional security risks have also been reported in deployed systems that use zero-knowledge proof protocols, including Zcash \cite{zcash_counterfeiting}, zkSync Era \cite{chainlight_zksync_2023}, and Polygon zkEVM \cite{verichains_polygon_zkevm_2024}. These cases show that unsafe implementation and deployment of zero-knowledge proof systems can undermine their intended security guarantees. It is therefore crucial to develop effective methods for analyzing and verifying the correctness of real-world zero-knowledge proof implementations.

Given these risks, evaluating the security and correctness of zero-knowledge proof systems is necessary to maximize security and achieve their intended guarantees. However, manual source-code auditing is costly, labor-intensive, and often ineffective. Various techniques have been proposed to automatically detect vulnerabilities in software and cryptographic implementations, such as fuzzing test. In addition, several tools have been specifically designed for zero-knowledge proof systems, including Circomspect \cite{trailofbits2023circomspect}, Picus \cite{pailoor2023automated}, SNARKProbe \cite{fan2024snarkprobe}, ZKFUZZ \cite{takahashi2025zkfuzz}, and OWBB23 \cite{ozdemir2025bounded}. However, most of these tools focus on detecting security issues in circuits rather than in zero-knowledge proof implementations more broadly \cite{chaliasos2024sok}. Furthermore, the tools that analyze vulnerabilities in the proving and verification phases typically require access to the target application's source code, which limits their practical applicability. This leaves an important gap: existing approaches are less suitable for deployed zero-knowledge proof applications when source code is unavailable, while real-world security evaluation increasingly requires techniques that can operate directly on binaries.

This limitation is particularly important in practice because many cryptographic libraries and applications, such as Apple’s iMessage and Microsoft’s BitLocker, are released as closed-source or proprietary software. As a result, independent researchers cannot easily inspect or audit the full implementation. Although many current zero-knowledge proof systems remain research-driven and are often released as open source, some real-world deployed zero-knowledge proof applications are not fully open source. For example, World ID \cite{world_advancing_decentralization} is a deployed zero-knowledge-proof-based identity system, and some components, such as its fraud-detection module, remain closed source. As zero-knowledge proof technology matures and is increasingly adopted in commercial products, it is reasonable to expect that more deployed zero-knowledge proof applications will include proprietary components, further increasing the difficulty of source-based security auditing.

Existing techniques for detecting security vulnerabilities in zero-knowledge proof applications and libraries are therefore insufficient for future real-world deployments, especially as more systems are integrated into proprietary or closed-source products. This motivates the development of new tools and techniques that can automatically detect security vulnerabilities in zero-knowledge proof applications without requiring source-code access. Such advances would reduce the cost of security evaluation and improve the practical applicability of vulnerability detection for future zero-knowledge proof systems.

%We developed a new tool based on the grey box fuzzing techquie, which does not require any source code access but have certain information about the potenial structure of targeting binary application to effectly. we also applt taint tracking techquie to access the fuzzing test and increase effectiocy. 

Formal verification provides strong guarantees, but it often requires precise formal specifications, substantial manual effort, and detailed access to the target implementation. These requirements make it difficult to apply at scale to real-world \zksnark applications, especially when source code is unavailable or the system includes many interacting components. Fuzzing, in contrast, is more automated and practical for deployed software because it can directly test implementations with concrete inputs and uncover unexpected logic errors without requiring a complete formal model. Traditional fuzzing is useful for detecting crashes and memory errors~\cite{AFLplusplus,google_afl_based_fuzzers_overview}. However, it is not sufficient for finding cryptographic logic errors. A cryptographic program may return a normal verification result even when that result is incorrect. Therefore, testing must compare the program’s behavior against the intended mathematical statement rather than merely observe whether the program crashes. Differential fuzzing provides a promising direction because it compares an implementation against a trusted reference. However, applying differential fuzzing to zero-knowledge proof software is difficult because each application may prove a different statement, and valid inputs must satisfy strict structural and mathematical requirements.

To address these challenges, we present \greyboxfuzz, a grey-box differential fuzzing framework for detecting cryptographic logic errors in \zksnark binary programs. \zksnark is a widely used class of zero-knowledge proof systems because it provides succinct proofs and fast verification, making it practical for real-world applications. Greybox fuzzing is particularly well suited to this setting because it can obtain limited internal information from the target application, such as execution coverage and high-level structure, without requiring access to the source code. This is important in practice because many deployed \zksnark applications are available only as binaries, making source-based analysis infeasible. Given binary proof-generation and proof-verification programs, the required cryptographic inputs, and a user-defined statement, \greyboxfuzz generates structured test inputs, executes the target binaries, compares their behavior against a ground-truth statement implementation, and reports inconsistencies.

% To address these challenges, we present \greyboxfuzz, a grey-box differential fuzzing framework for detecting cryptographic logic errors in zero-knowledge proof binary programs. Greybox fuzzing is a fuzzing approach that can obtain limited internal information from the target application, such as execution coverage, without requiring access to the source code, thereby improving testing efficiency while remaining practical to deploy. \greyboxfuzz targets a practical setting in which source code may not be available. Given binary proof-generation and proof-verification programs, the required cryptographic inputs, and a user-defined statement, \greyboxfuzz generates structured test inputs, executes the target binaries, compares their behavior against a ground-truth statement implementation, and reports inconsistencies.

We evaluate \greyboxfuzz on real-world \zksnark programs and known vulnerability patterns, including missing range checks, missing output constraints, under-constrained circuits, nondeterministic nullifiers, and other errors that may not cause crashes. We also evaluate the effectiveness of our taint-tracking-based branch identification by comparing automatically identified cryptographic-related branches with manually labeled branches in binary programs. Our results show that focusing fuzzing effort on cryptographic-related branches improves the efficiency of detecting cryptographic logic errors compared with general differential fuzzing.

In summary, this paper makes the following contributions:

\begin{itemize}
    \item We introduce \greyboxfuzz, a grey-box differential fuzzing framework for detecting cryptographic logic errors in \zksnark binary programs without requiring source-code access.

    \item We implement \greyboxfuzz as a practical tool that can be used to test \zksnark applications without requiring users to have expert knowledge of zero-knowledge proofs or cryptography.

    \item We evaluate the performance of \greyboxfuzz against prior techniques in related domains and show that it achieves practical performance for real-world \zksnark applications.

    \item \greyboxfuzz successfully detects security vulnerabilities and bugs in seven \zksnark applications drawn from both existing libraries and academic works.
    \item To the best of our knowledge, we are the first to focus on binary \zksnark application security evaluation and vulnerabilities detection.

\end{itemize}

% \begin{itemize}
%     \item We introduce a grey-box differential fuzzing framework for detecting cryptographic logic errors in zero-knowledge proof binary programs without requiring source-code access.

%     \item We design a statement-based testing oracle that uses a high-level Charm specification as the ground truth for differential fuzzing.

%     \item We develop a structure-aware input formatter that generates and mutates cryptographic inputs according to their type structure, mathematical constraints, and statement-level validity.

%     \item We propose a cryptographic-branch-focused coverage monitor that combines static taint tracking and dynamic tracing to guide fuzzing toward security-relevant computation.

%     \item We build an error locator that helps identify whether a detected inconsistency originates from the circuit, input conversion, proof generation, or proof verification.
% \end{itemize}
